\documentclass[11pt]{article}
\usepackage[T1]{fontenc}
\usepackage[utf8]{inputenc}
\usepackage{amsmath,amssymb,amsthm,mathtools}
\usepackage{geometry,graphicx,booktabs,longtable,tabularx,array}
\usepackage{enumitem,needspace,xcolor,tikz,tcolorbox}
\usetikzlibrary{arrows.meta,positioning,calc,fit,backgrounds,decorations.pathreplacing,matrix,patterns}
\tcbuselibrary{breakable}
\usepackage{caption,hyperref,bookmark,cleveref,microtype}
\usepackage{komorebi-publication}
\renewcommand{\KomorebiTitle}{Komorebi Space}
\renewcommand{\KomorebiSubtitle}{Coherent realizations, obstructions, and descent}
\renewcommand{\KomorebiSeries}{Foundations of the Realization Calculus}
\renewcommand{\KomorebiKind}{Monograph}
\renewcommand{\KomorebiEdition}{Budget continuation edition 1.10.0 / 12 October 2026}
\renewcommand{\KomorebiShortTitle}{Komorebi Space}
\renewcommand{\KomorebiProfile}{monograph}
\KomorebiSetup
\setcounter{tocdepth}{1}
\setcounter{secnumdepth}{2}
\setlist{nosep,leftmargin=1.7em,topsep=4pt}
\hypersetup{pdfkeywords={Komorebi Space,realization calculus,descent,obstruction,Lean,holonomy,closed range}}
\newtheorem{theorem}{Theorem}[section]
\newtheorem{proposition}[theorem]{Proposition}
\newtheorem{lemma}[theorem]{Lemma}
\newtheorem{corollary}[theorem]{Corollary}
\theoremstyle{definition}
\newtheorem{definition}[theorem]{Definition}
\newtheorem{example}[theorem]{Example}
\theoremstyle{remark}
\newtheorem{remark}[theorem]{Remark}
\numberwithin{equation}{section}
\newcommand{\R}{\mathbb R}
\newcommand{\Z}{\mathbb Z}
\newcommand{\Q}{\mathbb Q}
\newcommand{\N}{\mathbb N}
\newcommand{\K}{\mathbb K}
\newcommand{\V}{\mathcal V}
\newcommand{\W}{\mathcal W}
\newcommand{\id}{\operatorname{id}}
\newcommand{\im}{\operatorname{im}}
\newcommand{\ran}{\operatorname{ran}}
\newcommand{\Hom}{\operatorname{Hom}}
\newcommand{\Ob}{\operatorname{Ob}}
\newcommand{\Ext}{\operatorname{Ext}}
\newcommand{\Match}{\operatorname{Match}}
\newcommand{\Exact}{\operatorname{Exact}}
\newcommand{\dist}{\operatorname{dist}}
\newcommand{\Aut}{\operatorname{Aut}}
\newcommand{\Stab}{\operatorname{Stab}}
\newcommand{\Set}{\mathbf{Set}}
\newcommand{\Cat}{\mathbf{Cat}}
\newcommand{\norm}[1]{\left\lVert#1\right\rVert}
\newcommand{\abs}[1]{\left\lvert#1\right\rvert}
% Start the callout paragraph before anchors/color whatsits: in vertical mode
% those nodes would create a fresh break at the following paragraph glue.
\newcommand{\source}[1]{\par\nopagebreak[4]\noindent\phantomsection\label{src:#1}{\footnotesize\sffamily\color{KomorebiMuted}Source concordance: \hyperlink{claim:#1}{#1}.}\par\smallskip}
\newcommand{\code}[1]{\texttt{\detokenize{#1}}}
\newcommand{\chap}[1]{\par\Needspace{10\baselineskip}\section{#1}}
\newcommand{\partnote}[2]{\clearpage\part{#1}\noindent\begin{minipage}{.92\linewidth}\large #2\end{minipage}\par\vspace{12pt}}
\newtcolorbox{insight}{breakable,colback=KomorebiSoft!55,colframe=KomorebiPine!50,boxrule=.45pt,arc=1pt,left=9pt,right=9pt,top=7pt,bottom=7pt}
\newtcolorbox{scopebox}{breakable,colback=KomorebiWarm,colframe=KomorebiOchre!55,boxrule=.4pt,arc=1pt,left=9pt,right=9pt,top=7pt,bottom=7pt}
\usepackage{float}
\newcommand{\fixedplate}[3]{\begin{figure}[H]\centering\resizebox{.96\linewidth}{!}{\input{figures/#1}}\caption{#2}\label{#3}\end{figure}}
\newcommand{\plate}[3]{\begin{figure}[!htbp]\centering\resizebox{.96\linewidth}{!}{\input{figures/#1}}\caption{#2}\label{#3}\end{figure}}
\tikzset{box/.style={draw=KomorebiPine!70,rounded corners=2pt,fill=KomorebiSoft!50,align=center,inner sep=8pt,font=\small},arr/.style={-{Latex[length=2.2mm]},draw=KomorebiPine,thick},thinarr/.style={-{Latex[length=1.8mm]},draw=KomorebiMuted},lab/.style={font=\small,align=center},dot/.style={circle,fill=KomorebiPine,inner sep=2.3pt}}
\renewcommand{\arraystretch}{1.18}
\raggedbottom
\widowpenalty=7000\clubpenalty=7000
\setlength{\emergencystretch}{3em}

% Compact only the contents entries, without reducing body type size.
\usepackage{etoolbox}
\makeatletter
% Suppress the break inserted by endtrivlist after the final line only.
% Interior paragraphs and lines of long proofs remain normally breakable.
\newcommand{\KomorebiKeepSourceEnding}[1]{%
  \AtEndEnvironment{#1}{\@endparpenalty\@M}}
\forcsvlist{\KomorebiKeepSourceEnding}{proof,theorem,proposition,lemma,corollary,definition,example,remark}
\patchcmd{\l@section}{1.0em}{0.6em}{}{}
\makeatother

% Historical statistics and current budget-continuation evidence counts.
% Historical counts retain the scope of their original accepted records.
% Cumulative v1.8.0 counts below are bound to continuation/VERIFICATION_SUMMARY.json.
\newcommand{\EditionInexactChecks}{18,728}
\newcommand{\EditionFocusedTargets}{7}
\newcommand{\EditionFocusedDeclarations}{223}
\newcommand{\EditionFocusedNewDeclarations}{97}
\newcommand{\EditionFocusedTheorems}{51}
\newcommand{\EditionFocusedReferences}{30}
\newcommand{\EditionClaims}{80}
\newcommand{\EditionReferences}{211}
% Written-only KS entries. The current WR count is EditionWRWrittenOnly below.
\newcommand{\EditionWrittenOnly}{0}
\newcommand{\EditionOriginalChecks}{14,224}
\newcommand{\EditionRecoveredChecks}{2,767}
\newcommand{\EditionRevisionChecks}{17,213}
\newcommand{\EditionExpansionChecks}{96,345}
\newcommand{\EditionDeepeningChecks}{13,817}
\newcommand{\EditionTotalChecks}{197,482}
\newcommand{\EditionRegularizationChecks}{34,388}
\newcommand{\EditionWrittenExtensions}{16}
% Historical v1.2.0 canonical core counts, not the current cumulative totals.
\newcommand{\EditionFormalModules}{138}
\newcommand{\EditionAuditedDeclarations}{4,745}
\newcommand{\EditionNewTheorems}{54}
\newcommand{\EditionFormalTheorems}{1,507}
\newcommand{\EditionNewConstructions}{17}
\newcommand{\EditionReplays}{1}
\newcommand{\EditionNegativeControls}{7}
% Accepted v1.8.0 cumulative baseline and subsequent WR01--WR02 continuation.
\newcommand{\EditionCumulativeBaseModules}{141}
\newcommand{\EditionCumulativeBaseTargets}{144}
\newcommand{\EditionCumulativeBaseDeclarations}{4,842}
\newcommand{\EditionContinuationModules}{3}
\newcommand{\EditionContinuationTargets}{4}
\newcommand{\EditionContinuationTheorems}{22}
\newcommand{\EditionContinuationDeclarations}{30}
\newcommand{\EditionCumulativeModules}{144}
\newcommand{\EditionCumulativeRoots}{3}
\newcommand{\EditionCumulativeTargets}{147}
\newcommand{\EditionCumulativeDeclarations}{4,872}
\newcommand{\EditionCumulativeTheorems}{1,580}
\newcommand{\EditionCumulativePositiveFiles}{4}
\newcommand{\EditionCumulativeCompilerNegativeFiles}{14}
\newcommand{\EditionCumulativePolicyNegativeFiles}{2}
\newcommand{\EditionWRComplete}{9}
\newcommand{\EditionWRPartial}{2}
\newcommand{\EditionWRWrittenOnly}{5}
% Version 1.10.0 budget continuation. Historical counters above remain unchanged.
\newcommand{\EditionBudgetModules}{148}
\newcommand{\EditionBudgetRoots}{3}
\newcommand{\EditionBudgetTargets}{151}
\newcommand{\EditionBudgetNewModules}{4}
\newcommand{\EditionBudgetTheorems}{1,619}
\newcommand{\EditionBudgetNewTheorems}{39}
\newcommand{\EditionBudgetControls}{30}
\newcommand{\EditionBudgetPositiveFiles}{8}
\newcommand{\EditionBudgetCompilerNegativeFiles}{20}
\newcommand{\EditionBudgetPolicyNegativeFiles}{2}
\newcommand{\EditionBudgetDeclarations}{4,935}
\newcommand{\EditionBudgetNewDeclarations}{63}


% Keep readable half-page figures with their surrounding exposition.
\renewcommand{\topfraction}{.88}
\renewcommand{\bottomfraction}{.80}
\renewcommand{\textfraction}{.10}
\renewcommand{\floatpagefraction}{.78}
\setcounter{topnumber}{2}
\setcounter{bottomnumber}{2}
\setcounter{totalnumber}{3}

% WR evidence is synchronized with the budget continuation.
% WR evidence labels synchronized with the version 1.10.0 budget continuation.
\expandafter\def\csname WRstatusWR01\endcsname{Lean counterpart verified, including preservation of the same correction budget.}
\expandafter\def\csname WRstatusWR02\endcsname{Lean counterpart verified for nested nonempty closed affine constraints.}
\expandafter\def\csname WRstatusWR03\endcsname{Lean counterpart verified for attainment, profile geometry, the limit floor, and exact feasibility.}
\expandafter\def\csname WRstatusWR04\endcsname{Lean counterpart verified for the attained sharp dual maximum at every nonnegative budget.}
\expandafter\def\csname WRstatusWR05\endcsname{Lean counterpart verified for profile uncertainty and nominal dual lower certificates on the same correction space.}
\expandafter\def\csname WRstatusWR06\endcsname{Lean counterpart verified in the cumulative scope. Examples and prose correspondence remain separately scoped.}
\expandafter\def\csname WRstatusWR07\endcsname{Lean counterpart verified in the cumulative scope. Examples and prose correspondence remain separately scoped.}
\expandafter\def\csname WRstatusWR08\endcsname{Complete written proof. No new Lean counterpart is assigned.}
\expandafter\def\csname WRstatusWR09\endcsname{Partial Lean coverage. Pointwise estimates are verified. The limiting or additional conclusions retain written proofs.}
\expandafter\def\csname WRstatusWR10\endcsname{Complete written proof. No new Lean counterpart is assigned.}
\expandafter\def\csname WRstatusWR11\endcsname{Complete written proof. No new Lean counterpart is assigned.}
\expandafter\def\csname WRstatusWR12\endcsname{Complete written proof. No new Lean counterpart is assigned.}
\expandafter\def\csname WRstatusWR13\endcsname{Lean counterpart verified in the cumulative scope. Examples and prose correspondence remain separately scoped.}
\expandafter\def\csname WRstatusWR14\endcsname{Lean counterpart verified in the cumulative scope. Examples and prose correspondence remain separately scoped.}
\expandafter\def\csname WRstatusWR15\endcsname{Complete written proof. No new Lean counterpart is assigned.}
\expandafter\def\csname WRstatusWR16\endcsname{Partial Lean coverage. Pointwise estimates are verified. The limiting or additional conclusions retain written proofs.}

\newcommand{\written}[1]{\par\nopagebreak[4]\smallskip{\footnotesize\sffamily\color{KomorebiMuted}Extension \hyperref[app:written]{#1}. \csname WRstatus#1\endcsname}\par\smallskip}

\begin{document}
\KomorebiCover
\pagenumbering{roman}
\section*{About this edition}
This monograph develops Komorebi Space as a theory of coherent realizations relative to a specified system of mathematical views. Its central constructions are ordinary limits, relative extension spaces, obstruction quotients, cocycle groupoids, an explicit stack on a composition site, and quantitative repair operators. The point of the synthesis is to keep these constructions connected without erasing the differences between their hypotheses.

Definitions precede existence questions. Existence, uniqueness, and numerical stability remain separate. Changes of representation retain the distinction between the new representation and the data it must preserve.

Publication version \textbf{1.10.0} adds formal counterparts for the three budget theorems WR03--WR05. It retains the expanded proofs from version 1.7.0 and the full proofreading pass from version 1.9.0. This edition also corrects the distinction between a nearest reachable image and its residual. The analytic appendix supplies the projection, adjoint, Banach inverse, and Neumann-series tools used in the arguments.

The eighty source-linked KS entries retain their declaration correspondences. The enlarged project has \EditionBudgetModules\ production modules and \EditionBudgetRoots\ roots. The new budget statements preserve unrestricted real Hilbert spaces and include zero budget. Appendices~\ref{app:concordance} and~\ref{app:written} separate the historical build records from the current verification and identify the public declarations.

Of the sixteen WR groups, WR01--WR07, WR13, and WR14 have complete typed counterparts for their principal statements. WR09 and WR16 have formal pointwise estimates but written limiting consequences. The other five groups remain written-only. The printed weak-subsequence lemma remains a written argument. The formal WR01 proof reaches the same bounded realization conclusion by a separate compactness route. Appendix~\ref{app:written} gives the exact coverage.

The budget continuation supplies new Lean compilation and focused proof review. Historical finite Python suites are preserved and were not rerun. Publication checks cover the changed prose, manuscript structure, references, the PDF build, and rendered pages. This edition does not claim a fresh theorem-by-theorem referee audit of the entire monograph.

Lean's compiler, kernel, and pinned external dependency cache remain trusted. This is not an independent rebuild of mathlib or a proof that a Python program implements every formal construction. Source correspondence identifies the actual declarations and hypotheses. It does not turn explanatory prose into a kernel-checked language.

All mathematical figures are reproducible vector diagrams or Python plots of explicit formulas. The package contains their sources, the manuscript, exact checks, evidence, and progress ledger. The retained visual and Lean guides provide companion routes through the examples and historical evidence. This monograph states the updated formal coverage.

\begin{scopebox}
\textbf{Reading scope.} The word \emph{space} does not assert one universal topological object containing all of mathematics. The strict construction is a set of coherent families. Symmetry is retained in a groupoid, and effective descent is proved for a specified site. Analytic estimates apply only after a normed or Hilbert model has been declared. These are related layers, not interchangeable definitions.
\end{scopebox}

\clearpage
\section*{Orientation}
A description at one scale often omits information visible at another. A formula in one representation may fail to commute with a change of representation. Local pieces may each be valid and still fail to assemble. Even when assembly is possible, there may be many completions, and small discrepancies may or may not imply proximity to an exact answer.

Komorebi Space treats these questions in a common order. First specify the views and the maps between them. Then define the coherent realizations. Next identify the information that must remain fixed. Finally determine the obstruction to extension, the freedom in a successful extension, and the conditions under which repair is stable.

The resulting program has three central distinctions:
\[
\begin{gathered}
\text{a collection of appearances}\quad\ne\quad\text{a coherent realization},\\
\text{existence of an exact answer}\quad\ne\quad\text{a stable repair estimate},\\
\text{one isomorphism class}\quad\ne\quad\text{absence of symmetries}.
\end{gathered}
\]
The monograph makes these distinctions into explicit constructions and theorems.

\subsection*{The route through the book}
Part I constructs coherent sections and follows their obstructions as more data are prescribed. Part II develops resolution towers, finite observations, split repair coordinates, and multiscale reconstruction. Part III retains symmetries, proves effective descent, and distinguishes coherent phases from removable ones. Part IV develops analytic stability and an independent geometry of transformation costs. Part V follows approximate answers to their limits, distinguishes a coherent point from its intended Hilbert realization, and then derives the bounded-budget and regularization criteria that diagnose the difference. Part VI works through two finite loop models and returns to mathematical semantics. Appendix~\ref{app:workshop} contains fifteen solved problems. Appendix~\ref{app:hypotheses} gives a hypothesis atlas, and Appendix~\ref{app:written} records the exact evidence scope of the extensions. Appendix~\ref{app:proof-tools} proves the analytic background used in their arguments.

A short technical route is Sections~\ref{sec:views}, \ref{sec:obstruction}, \ref{sec:stack}, \ref{sec:closed}, and~\ref{sec:motion}. A reader interested in the role of packets and residuals should include Sections~\ref{sec:split} and~\ref{sec:cascade}. The graph models in Section~\ref{sec:graph} explain how a categorical realization problem acquires an actual numerical condition number. Sections~\ref{sec:observation} and~\ref{sec:repair-transport} give the representation-theoretic links. Sections~\ref{sec:signal-loop}, \ref{sec:twisted-loop}, and~\ref{sec:escape} can also be read first as concrete motivations. The last of these connects finite realization, closure, and failure of protected infinite existence in one calculation. Sections~\ref{sec:bounded-realization} through~\ref{sec:diagonal-regularization} give its positive converse, fixed-budget diagnostic, and regularized continuation. Sections~\ref{sec:inexact-repair} and~\ref{sec:stopping-rules} then explain how to certify an inexact calculation and choose its stopping tolerance.

\subsection*{What is being contributed}
The limit universal property, gauge action groupoids, standard stack definitions, orthogonal projection, and bounded inverse principles belong to established mathematics. They are used here as such. The contribution of this presentation is their explicit integration with the project's actual view systems, protected correction spaces, resolution differentials, graph transports, and checked examples. Historical novelty of every component is not inferred from formal verification. Conversely, the fact that the ingredients are classical does not remove the need to prove the comparison maps and to retain the right hypotheses.

\clearpage\begingroup\setlength{\parskip}{0pt}\tableofcontents\endgroup
\clearpage\pagenumbering{arabic}
\input{chapters/01_views}
\input{chapters/02_obstructions}
\input{chapters/02a_prescription}
\input{chapters/03_towers}
\input{chapters/04_symmetry}
\input{chapters/04a_boundary_lab}
\input{chapters/04b_phase}
\addtocontents{toc}{\protect\newpage}
\input{chapters/05_geometry}
\input{chapters/05a_certificates}
\input{chapters/05b_transport_cost}
\input{chapters/05c_iteration}
\input{chapters/05d_floors}
\input{chapters/05e_escape}
\clearpage
\chap{A bounded repair budget can recover existence}\label{sec:bounded-realization}
The preceding example does not merely lack a convergent iteration. It has no bounded family of increasingly accurate protected corrections. This suggests a converse question. If approximate corrections stay in one bounded region, does an exact answer exist?

For bounded linear equations in a Hilbert space, the answer is yes. This is an existence statement, not a claim that the given sequence converges in norm. A stronger conclusion becomes available when the candidates are the least norm solutions of nested finite constraint systems.

\begin{scopebox}
This chapter and the next three retain the WR01--WR12 identifiers. Every proof is printed. WR01 and WR02 now have complete formal counterparts in the accepted version 1.8.0 combined Lean project. The printed proof of WR01 uses the weak-subsequence lemma below, which remains a written background argument. Acceptance of the formal counterpart does not by itself verify that printed proof route. The later regularized optimizer and parameter results WR06--WR07 are also included in the accepted combined project. Fixed-parameter estimates within WR09 are formalized, but its limit statement is not. Version 1.10.0 adds complete formal counterparts for WR03--WR05: the attained budget profile, its sharp dual certificate, and its robustness on a fixed correction space. Each result carries its own evidence label. Exact finite models have separate executable checks. The Hilbert and regularization background is classical \cite{clason-inverse}. No novelty is claimed for weak compactness or quadratic regularization.
\end{scopebox}

\subsection{The correction space is the protected kernel}
Let $E,Y,Z$ be real Hilbert spaces, let $d:E\to Y$ and $p:E\to Z$ be bounded linear maps, and fix a raw candidate $f\in E$. Write
\[
 H=\ker p,\qquad T=d|_H:H\to Y,\qquad r=df.
\]
The space $H$ is a closed Hilbert subspace with the norm inherited from $E$. An allowed correction is $a\in H$, and its repaired candidate is $f-a$. Thus
\[
 p(f-a)=pf,\qquad d(f-a)=r-Ta.
\]
The protected exact problem is precisely $Ta=r$. Restricting the domain before optimizing is important. A penalty on $\norm{pa}$ would only discourage a change in the protected data, not prohibit it.

The adjoint below always means $T^*:Y\to H$, taken with the inherited Hilbert metric on $H$. In ambient coordinates this is the projection of $d^*$ onto $H$. A coordinate basis of $\ker p$ need not be orthonormal. Its Gram matrix must then be included in any calculation of correction energy.

\subsection{The compactness needed for the existence argument}
\begin{lemma}[Bounded sequences have weakly convergent subsequences]\label{lem:weak-subsequence}
Every bounded sequence $(a_n)$ in a real Hilbert space has a subsequence converging weakly to some $a$. Moreover,
\[
 \norm a\le\liminf_k\norm{a_{n_k}}.
\]
\end{lemma}
\begin{proof}
Choose $M\ge0$ with $\norm{a_n}\le M$. If $M=0$, every term is zero and the conclusion is immediate. Otherwise let $H_0$ be the closed linear span of the countably many vectors $a_n$. Finite rational linear combinations of those vectors form a countable dense set, so $H_0$ is separable. Apply Gram--Schmidt to a countable dense sequence, omitting zero remainders. This gives either a finite orthonormal basis or a countable orthonormal sequence $(e_j)$ whose closed span is $H_0$. The ambient Hilbert space itself need not be separable.

Each scalar sequence $\langle a_n,e_j\rangle$ is bounded by $M$. For the first coordinate choose a convergent subsequence, for the second choose a subsequence of that subsequence, and continue. Taking the $k$th term from the $k$th nested subsequence, with increasing original indices, yields $a_{n_k}$ for which every fixed coordinate converges. Call the coordinate limits $c_j$. In the finite-dimensional case only finitely many successive selections are required.

For each finite $J$, Bessel's inequality and passage to the limit give
\[
 \sum_{j=1}^Jc_j^2
 =\lim_k\sum_{j=1}^J\langle a_{n_k},e_j\rangle^2\le M^2.
\]
The nondecreasing partial sums are therefore bounded. Orthogonality makes the partial vector sums $\sum_{j\le J}c_je_j$ Cauchy, so completeness defines $a\in H_0$ with $\norm a\le M$.

To verify weak convergence against $h\in H_0$, choose a finite basis combination $h_J$ close to $h$. Then
\[
 |\langle a_{n_k}-a,h\rangle|
 \le |\langle a_{n_k}-a,h_J\rangle|
      +2M\norm{h-h_J}.
\]
For a prescribed tolerance, choose $h_J$ first to control the second term and then $k$ to control the first by finitely many coordinate limits. For a general $h$ in the ambient space, write $h=P_{H_0}h+h_\perp$. The orthogonal component pairs to zero with every $a_{n_k}$ and with $a$, reducing to the preceding case. This proves weak convergence against every test vector.

If $a\ne0$, testing against $a/\norm a$ gives
\[
 \norm a=\lim_k\langle a_{n_k},a/\norm a\rangle
 \le\liminf_k\norm{a_{n_k}}.
\]
The zero limit case follows from nonnegativity. Thus the lower-semicontinuity conclusion concerns the actual norms on the selected subsequence, not only the original bound $M$.
\end{proof}

\begin{theorem}[Bounded approximate feasibility]\label{WR01}
For a bounded linear map $T:H\to Y$ between real Hilbert spaces and $r\in Y$, the following are equivalent:
\[
 \begin{split}
 &r\in\ran T,\\
 &\text{there is a bounded sequence }(a_n)\subset H
       \text{ with }\norm{Ta_n-r}\longrightarrow0.
 \end{split}
\]
If all $\norm{a_n}\le M$, an exact solution exists with norm at most $M$.
\end{theorem}
\begin{proof}
If $Ta=r$, the constant sequence $a_n=a$ is bounded and has zero residual. For the converse, let $\norm{a_n}\le M$ and $Ta_n\to r$ in norm. By the preceding lemma there is a subsequence $a_{n_k}$ converging weakly to $a\in H$ with $\norm a\le M$.

For every fixed $y\in Y$, the adjoint identity and weak convergence give
\[
 \langle Ta,y\rangle
 =\langle a,T^*y\rangle
 =\lim_k\langle a_{n_k},T^*y\rangle
 =\lim_k\langle Ta_{n_k},y\rangle.
\]
The last limit equals $\langle r,y\rangle$ because norm convergence implies convergence of each continuous pairing. Thus $\langle Ta-r,y\rangle=0$ for every $y$. Taking $y=Ta-r$ gives $\norm{Ta-r}^2=0$, hence $Ta=r$.

This proof does not claim that $Ta_{n_k}$ converges in norm merely because $a_{n_k}$ converges weakly. Its norm convergence to $r$ was a separate hypothesis. The boundedness of the candidate corrections is used to obtain a weakly convergent subsequence, and their accurate residuals are used to identify its image exactly. The bound $\norm a\le M$ comes from weak lower semicontinuity, so the exact answer stays inside the same closed budget ball.
\end{proof}
\written{WR01}

No closed range assumption appears. The theorem concerns a single datum and a bounded family of its corrections. Closed range is the stronger, uniform statement that every attainable datum admits a correction controlled by one residual gain.

An immediate contrapositive is useful. If $r\in\overline{\ran T}\setminus\ran T$, then every sequence with $Ta_n\to r$ satisfies $\norm{a_n}\to\infty$. Otherwise infinitely many terms would lie in one bounded ball, contradicting the theorem along that subsequence. This recovers the escape mechanism of Section~\ref{sec:escape} without using the diagonal formula.

\subsection{Existence does not make an arbitrary sequence converge strongly}
Take $H=Y=\ell^2$ and $Te_n=e_n/n$. The vectors $a_n=e_n$ have norm one and satisfy $Ta_n\to0$. An exact answer exists, namely zero. Yet $\norm{e_n-e_m}=\sqrt2$ for $n\ne m$, so no subsequence converges in norm.

The conclusion of WR01 is the existence of a weak limit along a subsequence that solves the equation. It is not summable movement, a strong convergence rate, or a selection rule for the original candidates. The minimum correction for the zero datum is zero, which this particular sequence never approaches in norm.

\subsection{Minimum energy finite repairs do converge strongly}
\fixedplate{31_nested_energy}{Minimum energy finite corrections behave differently for two data under the same diagonal operator. Bounded partial energy produces a strong Hilbert realization. Escaping partial energy accompanies a coherent product-space point outside the Hilbert space. The attainable energy lies below the explicit bound two.}{fig:nested-energy}

Suppose now that $A_1\supseteq A_2\supseteq\cdots$ are nonempty closed affine subsets of a real Hilbert correction space $H$. A typical $A_n$ consists of corrections satisfying the first $n$ bounded linear observation equations. The sets are affine because the equations are linear, and closed because the observations are continuous.

Let $a_n$ be the unique point of $A_n$ nearest to zero, and put $m_n=\norm{a_n}$. The Hilbert projection theorem supplies these points. It can also be proved directly: a sequence approaching the minimum norm is Cauchy by the parallelogram identity, its limit belongs to the closed convex set, and the same identity gives uniqueness.

\begin{theorem}[A uniform finite budget yields a strong realization]\label{WR02}
For the nested affine sets above, $m_n$ is nondecreasing, and for $m\ge n$,
\begin{equation}\label{eq:nested-energy}
 \norm{a_m-a_n}^2=m_m^2-m_n^2.
\end{equation}
Their intersection is nonempty if and only if $\sup_n m_n<\infty$. In that case $a_n$ converges in norm to the unique least norm point $a_*\in\bigcap_n A_n$, with
\[
 \norm{a_*-a_n}^2=\norm{a_*}^2-\norm{a_n}^2.
\]
If the intersection is empty, then $m_n\to\infty$.
\end{theorem}
\begin{proof}
The nearest-point theorem for nonempty closed convex subsets of a Hilbert space, proved in Appendix~\ref{app:proof-tools}, supplies $a_n$. Write $A_n=a_n+L_n$, where $L_n$ is its closed direction subspace. For $\ell\in L_n$ and every real $t$, the vector $a_n+t\ell$ belongs to $A_n$. Minimality implies
\[
 0\le\norm{a_n+t\ell}^2-\norm{a_n}^2
 =2t\langle a_n,\ell\rangle+t^2\norm\ell^2.
\]
Allowing both signs of arbitrarily small $t$ forces $\langle a_n,\ell\rangle=0$.

If $m\ge n$, nesting gives $a_m\in A_n$, so $a_m-a_n\in L_n$. Expanding its norm therefore gives
\[
 \norm{a_m}^2
 =\norm{a_n+(a_m-a_n)}^2
 =\norm{a_n}^2+\norm{a_m-a_n}^2.
\]
This is the exact energy identity. It also makes $m_n^2$ nondecreasing.

Suppose these energies are bounded, and let $L=\lim_n m_n^2$. Given $\epsilon>0$, choose $N$ with $L-m_N^2<\epsilon^2$. For $m\ge n\ge N$,
\[
 \norm{a_m-a_n}^2=m_m^2-m_n^2\le L-m_N^2<\epsilon^2.
\]
Thus $(a_n)$ is Cauchy and has a norm limit $a_*$. For each fixed $n$, all terms from that index onward belong to the closed set $A_n$, so $a_*\in A_n$. This places $a_*$ in their entire intersection.

If $a$ lies in the intersection, then $\norm a\ge\norm{a_n}$ for every $n$. Norm continuity gives $\norm a\ge\norm{a_*}$. Uniqueness follows from the nearest-point theorem applied to the intersection, which is closed and convex and now known to be nonempty. Alternatively, the midpoint of two different minimizers would have strictly smaller norm by the parallelogram identity. Taking $m\to\infty$ in the finite energy identity yields
\[
 \norm{a_*-a_n}^2=\norm{a_*}^2-\norm{a_n}^2.
\]
Conversely, any element of the intersection bounds every minimum norm from above. Therefore an empty intersection forces the nondecreasing sequence $m_n$ to be unbounded. A nondecreasing unbounded nonnegative sequence tends to infinity, giving the last conclusion.
\end{proof}
\written{WR02}

The finite answers need not extend one another coordinate by coordinate. Their shared feature is minimum energy in one common ambient Hilbert space. This is the extra geometry that forces their differences to be orthogonal to the earlier minimizers.

For the attainable diagonal datum $r_n=1/n^2$ under $Tu=(u_n/n)$, with $n\ge1$, the first $N$ equations force correction coordinate $n$ to equal $1/n$ for $n\le N$. Their minimum energy completion is zero beyond $N$, and
\[
 m_N^2=\sum_{n=1}^{N}\frac1{n^2}\le2.
\]
These corrections converge strongly to $(1/n)$. For the escaping datum $r_n=1/n$, the corresponding energy is $m_N^2=N$. Every finite equation is again solved, but the budget test now rules out an ambient solution.


\begin{insight}
Finite consistency is not enough by itself. For nested closed affine problems in one Hilbert space, a uniform bound on the minimum correction norms is exactly what upgrades finite feasibility to ambient feasibility. This is a datum-specific test, not a uniform inverse bound for the operator.
\end{insight}

\chap{The attainable residual at a fixed correction budget}\label{sec:budget-frontier}
The unrestricted floor $\dist(r,\ran T)$ forgets how large a correction must be. To retain that information, fix a nonnegative budget $M$ and define
\[
 \beta_{T,r}(M)=\inf_{\norm a\le M}\norm{r-Ta}.
\]
The norm is the actual correction norm in $H=\ker p$. It is not the norm of arbitrary coordinates used to represent that kernel. The function $\beta$ describes the smallest residual compatible with a declared amount of allowed change.

\subsection{At every finite budget the residual minimum exists}
\begin{theorem}[Attainment and shape of the budget profile]\label{WR03}
For a bounded linear map between real Hilbert spaces, the infimum defining $\beta(M)$ is attained for every $M\ge0$. The profile is nonincreasing and convex. For $0\le M_1\le M_2$,
\[
 0\le\beta(M_1)-\beta(M_2)\le\norm T(M_2-M_1).
\]
Writing $C=\overline{\ran T}$, $b=P_Cr$, and $q=\norm{r-b}$, one has
\[
 \lim_{M\to\infty}\beta(M)=q.
\]
The equality $\beta(M)=q$ at a finite budget is equivalent to the existence of $a$ with $Ta=b$ and $\norm a\le M$. In particular,
\[
 \beta(M)=0\quad\Longleftrightarrow\quad
 \exists a\in H:\ Ta=r,\ \norm a\le M.
\]
\end{theorem}
\begin{proof}
The feasible ball is nonempty because it contains zero, and the infimum lies between zero and $\norm r$. Choose $a_n$ in that ball with $\norm{r-Ta_n}\to\beta(M)$. A weakly convergent subsequence has limit $a$ with $\norm a\le M$. Bounded linearity makes $r-Ta_{n_k}$ converge weakly to $r-Ta$. Applying weak lower semicontinuity in the target gives
\[
 \norm{r-Ta}\le\liminf_k\norm{r-Ta_{n_k}}=\beta(M).
\]
The reverse inequality holds by the definition of the infimum, so $a$ attains it. When $M=0$ the only feasible correction is zero, and the same conclusion reads $\beta(0)=\norm r$.

For convexity, take minimizers $a_1,a_2$ at budgets $M_1,M_2$ and $t\in[0,1]$. The candidate $a_t=ta_1+(1-t)a_2$ satisfies
\[
 \norm{a_t}\le tM_1+(1-t)M_2,
\]
and its residual is $t(r-Ta_1)+(1-t)(r-Ta_2)$. The triangle inequality therefore gives
\[
 \beta(tM_1+(1-t)M_2)\le t\beta(M_1)+(1-t)\beta(M_2).
\]
Monotonicity follows from inclusion of balls. If $0\le M_1\le M_2$ and $M_2>0$, choose a minimizer $a_2$ at $M_2$ and put $a_1'=(M_1/M_2)a_2$. Then $\norm{a_1'}\le M_1$ and $\norm{a_2-a_1'}\le M_2-M_1$, so
\[
 \beta(M_1)\le\norm{r-Ta_1'}
 \le\beta(M_2)+\norm T(M_2-M_1).
\]
Together with monotonicity this proves the Lipschitz estimate. The omitted case $M_2=0$ has $M_1=0$ and is equality.

Since $r-b$ is orthogonal to $C$ and $b-Ta\in C$, every correction satisfies
\begin{equation}\label{eq:range-decomposition}
 \norm{r-Ta}^2=q^2+\norm{b-Ta}^2.
\end{equation}
Thus $\beta(M)\ge q$. Given $\epsilon>0$, range density in $C$ supplies one correction $a$ with $\norm{b-Ta}<\epsilon$. For every $M\ge\norm a$ it follows that
\[
 q\le\beta(M)\le\sqrt{q^2+\epsilon^2}.
\]
First letting the budget increase and then letting $\epsilon\downarrow0$ proves the limit.

Finally, if $\beta(M)=q$, choose an attaining correction in that ball. The orthogonal identity forces $\norm{b-Ta}^2=0$, hence $Ta=b$. Conversely such a correction attains residual $q$. Setting $q=0$ proves the exact-feasibility criterion. Attainment at fixed budget is essential to the forward implication: a zero unrestricted infimum would not itself supply a preimage.
\end{proof}
\written{WR03}

Although $\ran T$ may not be closed, the image $T(\{a:\norm a\le M\})$ of each closed budget ball is closed and convex. To see closedness, take a convergent sequence in that image, lift it to the bounded ball, and apply the weak limit argument. The nonclosed range is the increasing union of these closed images. Passing to an unlimited budget can therefore lose closedness even though every fixed-budget optimization is well posed.

The nearest reachable image $y_M=P_{T(\{a:\norm a\le M\})}r$ is unique, so the best residual vector $r-y_M$ is unique as well. Corrections mapping to $y_M$ can differ by vectors in $\ker T$. Among those corrections there is a unique least norm correction. These two uniqueness statements should not be confused with injectivity of $T$.

\subsection{A dual witness can rule out every correction in the budget}
\begin{theorem}[A sharp budget certificate]\label{WR04}
For every $M\ge0$,
\begin{equation}\label{eq:budget-dual}
 \beta_{T,r}(M)=
 \max_{\norm z\le1}\bigl(\langle r,z\rangle-M\norm{T^*z}\bigr).
\end{equation}
Thus any single $z$ with $\norm z\le1$ and
\[
 \langle r,z\rangle-M\norm{T^*z}>\varepsilon
\]
certifies that no correction of norm at most $M$ has residual at most $\varepsilon$.
\end{theorem}
\begin{proof}
For every $\norm a\le M$ and $\norm z\le1$, the adjoint identity and Cauchy--Schwarz give
\[
\begin{aligned}
 \langle r,z\rangle-M\norm{T^*z}
 &=\langle r-Ta,z\rangle+\langle a,T^*z\rangle-M\norm{T^*z}\\
 &\le\langle r-Ta,z\rangle
 \le\norm{r-Ta}.
\end{aligned}
\]
Taking the infimum over corrections proves that every proposed dual value is at most $\beta(M)$. This part alone already proves the impossibility certificate.

To prove sharpness, let $a_M$ be an attaining correction from WR03 and put $v=r-Ta_M$. If $v=0$, then $\beta(M)=0$ and $z=0$ attains the dual bound. Assume henceforth $v\ne0$ and set $z=v/\norm v$.

Fix any feasible $a$ and put $h=a-a_M$. The segment $a_M+th$ is feasible for $0\le t\le1$. Minimality gives, for $t>0$,
\[
 0\le\norm{v-tTh}^2-\norm v^2
   =-2t\langle v,Th\rangle+t^2\norm{Th}^2.
\]
Divide by $t$ and let $t\downarrow0$. It follows that $\langle v,T(a-a_M)\rangle\le0$, hence
\[
 \langle T^*z,a\rangle\le\langle T^*z,a_M\rangle
 \quad(\norm a\le M).
\]
The support value of a Hilbert ball can be computed directly. If $T^*z\ne0$ and $M>0$, it is attained at $a=M T^*z/\norm{T^*z}$ and equals $M\norm{T^*z}$. If either factor is zero, every feasible pairing is zero at its maximum and the same formula holds. Thus
\[
 \langle T^*z,a_M\rangle=M\norm{T^*z}.
\]
Substituting this equality gives
\[
 \langle r,z\rangle-M\norm{T^*z}
 =\langle r-Ta_M,z\rangle=\norm v=\beta(M).
\]
We have constructed an attaining witness and thereby justified the word ``maximum.'' The proof uses one-sided variations in a convex ball, not an unproved exchange of an infimum and a supremum.
\end{proof}
\written{WR04}

In finite dimensions this is a checkable inequality, not merely an optimization prescription. A program may propose the correction or the witness. The displayed inequality can be verified independently of the proposal procedure. In a protected ambient calculation, use the restricted adjoint, not $d^*$ without projection onto $\ker p$.

\subsection{A two-dimensional calculation of the supporting witness}
Take
\[
 T=\begin{pmatrix}1&0\\0&1/2\end{pmatrix},\qquad r=(1,1),\qquad
 a=(1/2,2/5),\qquad v=r-Ta=(1/2,4/5).
\]
Then $T^*v=a$. At the budget $M=\norm a=\sqrt{41}/10$, the normalized vector $z=v/\norm v$ gives
\[
 \langle r,v\rangle-\norm a\norm{T^*v}
 =\frac{13}{10}-\frac{41}{100}=\frac{89}{100}=\norm v^2.
\]
Consequently $a$ is a globally optimal budgeted correction and $\beta(M)=\sqrt{89}/10$. It is not necessary to compare it with a dense numerical grid of competing vectors. The supporting witness proves optimality over the entire ball.

\fixedplate{32_budget_certificate}{The image of a Euclidean correction ball is an ellipse. The nearest reachable point has a residual normal to a supporting line. That same residual, normalized, is the dual witness in the exact budget formula. The plotted point and budget are the rational example in the text.}{fig:budget-certificate}

\subsection{How measurement and model uncertainty affect the budget}
\begin{theorem}[A robust budget profile]\label{WR05}
Suppose $\widetilde T:H\to Y$ and $\widetilde r\in Y$ satisfy
\[
 \norm{T-\widetilde T}\le\eta,\qquad
 \norm{r-\widetilde r}\le\delta,\qquad\eta,\delta\ge0.
\]
For every $M\ge0$,
\[
 \abs{\beta_{T,r}(M)-\beta_{\widetilde T,\widetilde r}(M)}
 \le\delta+\eta M.
\]
In particular, for $\norm z\le1$,
\[
 \beta_{T,r}(M)\ge
 \langle\widetilde r,z\rangle-M\norm{\widetilde T^*z}
       -\delta-\eta M.
\]
\end{theorem}
\begin{proof}
Fix a feasible correction $a$. The two residual vectors have difference
\[
 (r-Ta)-(\widetilde r-\widetilde Ta)
 =(r-\widetilde r)-(T-\widetilde T)a,
\]
whose norm is at most $c=\delta+\eta M$. The reverse triangle inequality therefore gives both inequalities
\[
 \norm{r-Ta}\le\norm{\widetilde r-\widetilde Ta}+c,
 \qquad
 \norm{\widetilde r-\widetilde Ta}\le\norm{r-Ta}+c.
\]
The domain of feasible corrections is identical in the two problems. Taking the infimum of each inequality over that common domain yields
\[
 \beta_{T,r}(M)\le\beta_{\widetilde T,\widetilde r}(M)+c,
 \qquad
 \beta_{\widetilde T,\widetilde r}(M)\le\beta_{T,r}(M)+c.
\]
These are precisely the absolute-value bound. One could equivalently evaluate at the two attained minimizers, but no relation between those minimizers is needed.

For any admissible $z$, WR04 for the nominal problem gives
\[
 \beta_{\widetilde T,\widetilde r}(M)
 \ge\langle\widetilde r,z\rangle-M\norm{\widetilde T^*z}.
\]
Combine this with the lower half of the preceding comparison. The result is a certificate against every correction in the original fixed ball, reduced by the explicit uncertainty allowance. If the correction space itself changed, the common-domain step would fail and this proof would no longer apply.
\end{proof}
\written{WR05}

This estimate identifies a practical boundary of increasingly large corrections. A model error $\eta$ is amplified by the correction budget even if the measured residual is tiny. The candidate-level inequality is
\[
 \norm{r-Ta}\le\norm{\widetilde r-\widetilde Ta}+\delta+\eta\norm a.
\]
The theorem compares maps on one fixed correction space. It does not cover an uncertain preservation map whose kernel itself changes.

\fixedplate{33_budget_uncertainty}{The nominal budget profile is surrounded by the rigorous uncertainty band $\delta+\eta M$. A large budget can improve the nominal fit while enlarging the uncertainty in that conclusion. The example uses $T=1$, $r=1$, $\delta=1/20$, and $\eta=1/10$.}{fig:budget-uncertainty}

\input{chapters/05h_regularized_repair}
\input{chapters/05i_diagonal_regularization}
\input{chapters/05j_inexact_certificates}
\input{chapters/05k_stopping_rules}
\input{chapters/06_cases}
\chap{Mathematical meaning and the remaining boundaries}\label{sec:meaning}
\subsection{A relative theory can be complete without being absolute}
The strict space \(\Gamma(F)\) is canonical relative to its diagram: its universal property determines it up to the appropriate unique isomorphism. The additive obstruction is canonical relative to its admissible corrections. A marked realization problem is invariant under the specified comparison equivalences. A repair can be canonical relative to a protected observable, or orthogonal relative to a fixed inner product.

These are meaningful and different kinds of canonicity. They need not culminate in one object independent of every choice of views, coefficients, topology, semantics, and geometry. Some changes of view intentionally forget information. Others strengthen the question. Treating every such change as an equivalence would remove the very distinctions the calculus is designed to preserve.

The appropriate central statement is therefore not that every possible mathematical presentation lives in one fixed universal space. It is that a declared system of representations can be organized into coherent realizations, with explicit extension, comparison, and repair problems. A proposed bridge between two such systems has to show what structure it preserves and what information it loses.

\subsection{Internal realization and intended objects}
Suppose a separate category \(\mathcal S\) contains the intended semantic objects of an application, and a realization construction gives a comparison functor
\[
 \mathcal R:\mathcal S\longrightarrow\mathcal K,
\]
where \(\mathcal K\) is the relevant category or groupoid of coherent data. Defining \(\mathcal K\) does not establish that \(\mathcal R\) is faithful, full, or essentially surjective. These are new mathematical questions about the actual semantic comparison.

Faithfulness asks whether distinct semantic arrows can become indistinguishable. Fullness asks whether every compatible comparison comes from a semantic one. Essential surjectivity asks whether every intended coherent datum is realized by an actual semantic object. Depending on the application, one may seek only a restricted image or a classified obstruction to belonging to it.

This distinction prevents a formal globalization of encoded data from being mistaken for an unrestricted arithmetic or physical existence theorem. The abstract construction does not manufacture missing semantic objects.

\subsection{Every finite test can pass while the infinite problem fails}
\begin{example}[Finite prefixes without a global point]\label{KS54}
Let \(X_n=\{m\in\N:m\ge n\}\), with bonding maps given by inclusion. Every level is nonempty, and every finite prefix has a coherent point. The inverse limit is nevertheless empty.
\end{example}
\begin{proof}
For a fixed terminal level \(N\), the integer \(N\) belongs to every fiber \(X_n\) with \(n\le N\). Assign that same integer at each such level. Every restriction is an inclusion, so it preserves the assigned integer and all finite coherence equations hold. This proves finite solvability with the quantifiers in the order ``for every \(N\), there exists a prefix.''

Suppose instead that \(x=(x_n)\) is a global point. Equality along each inclusion says that the underlying natural numbers of \(x_{n+1}\) and \(x_n\) are equal. Induction therefore makes every underlying coordinate equal to the number \(m\) underlying \(x_0\). Membership of \(x_{m+1}\) in \(X_{m+1}\) now gives \(m\ge m+1\), a contradiction. Thus no global point exists.

There is no conflict with a gluing theorem for prefixes. The prefix chosen above for length \(N+1\) restricts to the value \(N+1\), not to our chosen value \(N\) on the shorter prefix. Nor could another choice of one solution for every length give a matching family: its value at level zero would force the same forbidden global integer. Separate existence at every length is strictly weaker than existence of a compatible family of prefix solutions.
\end{proof}
\source{KS54}
This counterexample is formalized in \code{FinitePrefixBoundary.lean} using the original resolution and partial-section types. It explains why inverse-limit existence needs more than separate finite solvability. Finiteness, compactness, surjectivity, or another coherence principle must be proved in the relevant application. The example does not contradict the strict gluing theorem: a compatible matching family of all prefixes was never supplied.

\subsection{A positive infinite-existence criterion}
The failed example isolates a missing compactness or lifting hypothesis. Finite fibers supply a concrete positive criterion without requiring every bonding map to be surjective.

\begin{proposition}[Finite fibers and persistent images]\label{KS58}
Let \((X_n,\rho_n)\) be a sequential inverse system of nonempty finite sets. Its inverse limit is nonempty. More precisely, if \(\rho_{n,m}:X_m\to X_n\) is composite restriction, then
\[
 \im(\operatorname{ev}_n)=Y_n,
 \qquad Y_n=\bigcap_{m\ge n}\im\rho_{n,m}.
\]
Thus a prescribed value at level \(n\) extends precisely when it survives every later image.
\end{proposition}
\begin{proof}
Write \(I_{n,m}=\rho_{n,m}(X_m)\) for \(m\ge n\). Each \(I_{n,m}\) is nonempty. The composition law gives
\[
 I_{n,m+1}=\rho_{n,m}\bigl(\rho_m(X_{m+1})\bigr)\subseteq I_{n,m}.
\]
These are subsets of the finite set \(X_n\). Every strict decrease removes an element permanently, so only finitely many strict decreases are possible. There is an index \(s_n\ge n\) such that \(I_{n,m}=I_{n,s_n}\) for all \(m\ge s_n\). Consequently \(Y_n=I_{n,s_n}\) is nonempty.

We next prove that restriction induces a \emph{surjection} \(Y_{n+1}\to Y_n\), even though \(X_{n+1}\to X_n\) need not be surjective. For membership, let \(z\in Y_{n+1}\). Given \(m\ge n+1\), write \(z=\rho_{n+1,m}(w)\) for some \(w\in X_m\). Then \(\rho_n z=\rho_{n,m}(w)\). Together with the automatic condition at \(m=n\), this shows \(\rho_nz\in Y_n\).

For surjectivity, fix \(y\in Y_n\) and choose
\[
 m\ge\max\{s_n,s_{n+1},n+1\}.
\]
Since \(y\in I_{n,m}\), there is \(w\in X_m\) with \(\rho_{n,m}(w)=y\). Put \(z=\rho_{n+1,m}(w)\). Stabilization at level \(n+1\) gives \(z\in Y_{n+1}\), and composition gives \(\rho_nz=y\).

Choose \(x_0\in Y_0\) and recursively choose \(x_{n+1}\in Y_{n+1}\) above \(x_n\). The preceding surjectivity justifies each step. Adjacent coherence implies every composite coherence equation, so this gives a point. This is an existence construction using countably many choices, not an asserted effective procedure for locating the stabilization indices.

Finally, fix a level \(n\). Every global point satisfies \(x_n=\rho_{n,m}(x_m)\) for each \(m\ge n\), hence \(x_n\in Y_n\). Conversely, start with any \(y\in Y_n\), lift recursively above level \(n\) through the surjections just proved, and define the lower coordinates by \(x_i=\rho_{i,n}(y)\). The composite restriction laws verify coherence across and below level \(n\). The resulting point evaluates to \(y\). These two inclusions prove the asserted equality of evaluation images, not merely nonemptiness of the inverse limit.
\end{proof}
\source{KS58}
This is the standard finite-set inverse-limit principle, here with its sequential persistent-image proof~\cite{stacks-limits}. Its formalization in \code{FiniteRealizationExistence.lean} constructs an explicit equivalence with the pinned library's section type, then applies the finite cofiltered-system and eventual-range theorems. The printed sequential proof and this formal route establish the same stated criterion. Both use ordinary classical choice. Surjective bonding maps on arbitrary nonempty sets give a second sufficient condition by direct recursive lifting. Neither criterion says that arbitrary data in a non-surjective finite tower extend.

\subsection{A finite computation that exposes the preservation contract}
Consider a raw vector \(x=(2,0,5)\in\R^3\) with defect \(d(x)=x_1-x_2\). Unconstrained orthogonal repair gives \((1,1,5)\), with squared error two. If the first and third coordinates must be preserved, the unique exact repair is \((2,2,5)\), with squared error four.

Both outputs are exact. Only one satisfies the second preservation problem. Neither theorem contradicts the other because the admissible changes differ. An implementation that reports only zero residual would not distinguish the two tasks.

For a finite linear system over a field, the same question admits independently checkable positive and negative certificates. Let \(A\in\K^{m\times n}\) encode compatibility, \(P\in\K^{p\times n}\) the protected values, and \(x\in\K^n\) the raw candidate. Let the columns of \(N\) be a basis of \(\ker P\). The correction equation is
\[
 ANz=Ax,\qquad y=x-Nz.
\]
It is equivalent to the stacked system \(Ay=0\), \(Py=Px\).

\begin{proposition}[A finite certificate alternative]\label{KS57}
Exactly one of the following occurs:
\begin{enumerate}
\item There is \(y\) with \(Ay=0\) and \(Py=Px\). The full answer set is \(y+H\), where \(H=\ker A\cap\ker P\).
\item There are row vectors \(\lambda\in\K^{1\times m}\) and \(\mu\in\K^{1\times p}\) with
\[
 \lambda A=\mu P,\qquad\lambda Ax=1.
\]
\end{enumerate}
The negative witness annihilates every admissible correction defect, since \(\lambda AN=\mu PN=0\). The ambiguity dimension in the positive case is
\[
 \dim H=n-\operatorname{rank}\begin{pmatrix}A\\P\end{pmatrix}.
\]
\end{proposition}
\begin{proof}
Set \(T=\begin{pmatrix}A\\P\end{pmatrix}\) and \(b=\begin{pmatrix}0\\Px\end{pmatrix}\). A vector \(y\) is an exact protected answer if and only if \(Ty=b\). Perform Gaussian row operations on the augmented matrix, keeping their product \(U\). Each elementary row operation is invertible over the given field. Thus \(Ty=b\) and \(UTy=Ub\) have exactly the same solutions.

If no row of the reduced augmented system has zero left side and nonzero right side, assign arbitrary values to its free variables and solve for its pivot variables. This supplies a solution \(y_0\). For any vector \(y\),
\[
 Ty=b\quad\Longleftrightarrow\quad T(y-y_0)=0.
\]
Translation by \(y_0\) is therefore a bijection from \(\ker T\) to the answer set, with inverse \(y\mapsto y-y_0\). Rank--nullity applied to the actual linear map \(T:\K^n\to\K^{m+p}\) gives the ambiguity dimension. Dependent prescription rows do not alter this argument.

If the reduced system is inconsistent, choose a row with zero left side and right side \(\alpha\ne0\). Let \(w\) be the corresponding row of \(U\), and split it according to the two blocks as \(w=(u,v)\). The retained row operations give the literal identities
\[
 uA+vP=0,\qquad vPx=\alpha.
\]
Division is permitted because \(\alpha\) is a nonzero field element. Define \(\lambda=-\alpha^{-1}u\) and \(\mu=\alpha^{-1}v\). Multiplying the displayed equations by \(\alpha^{-1}\) gives
\[
 \lambda A=\mu P,\qquad \lambda Ax=\mu Px=1.
\]
For every admissible correction \(a\in\ker P\), these same row equations imply \(\lambda Aa=\mu Pa=0\). The witness therefore annihilates the entire correction image while taking value one on the required defect.

The alternatives cannot occur together. If \(y\) were an exact protected answer in the presence of such a witness, then
\[
 1=\lambda Ax=\mu Px=\mu Py=\lambda Ay=0.
\]
A field has \(1\ne0\), so this is impossible. Consistency or inconsistency of the reduced system proves exhaustiveness, and the contradiction proves exclusivity. Empty row or column blocks cause no extra division: normalization is used only in the inconsistent case, where the nonzero scalar has already been obtained. The proof does not transfer to integral coefficients by silently allowing the same division.
\end{proof}
\source{KS57}
The formal proof gives these literal row-vector equations over an arbitrary field. It uses algebraic dual separation of the stacked image rather than verifying the row-reduction program. The affine answer equivalence and rank-nullity formula are proved separately. Thus a normalized obstruction is an actual checkable witness, while universal correctness of the supplied Python elimination implementation remains a different question.

\subsection{One equation, three preservation problems}
For the following three cases work over \(\Q\), with \(A=(1,-1,0)\) and \(x=(2,0,5)\). With \(P(x)=(x_1,x_3)\), the unique answer is \((2,2,5)\). With \(P(x)=x_3\), the entire answer set is \(\{(t,t,5):t\in\Q\}\). With \(P(x)=(x_1,x_2)\), no answer exists. In this last case
\[
 \lambda=(1/2),\qquad\mu=(1/2,-1/2)
\]
is a normalized negative certificate over \(\Q\) or \(\R\): \(\lambda A=\mu P\) and \(\lambda Ax=1\).

The package implements this finite interface over \(\Q\) in \code{examples/protected_repair.py}. It returns a particular exact answer and an ambiguity basis, or a normalized dual witness. Its verifier checks the certificate equations directly. A Gaussian particular answer is \emph{not} advertised as the nearest answer of Theorem~\ref{KS56}. That is a separate geometric selection.

Exact field elimination does not solve the integer parity problem by changing its coefficient domain. Over \(\Q\), \(2m=1\) has the solution \(1/2\). Over \(\Z\), it has none. The rational interface rejects floating-point inputs rather than silently treating an approximate decimal as exact arithmetic. Its finite tests establish the tested implementation behavior, not an arbitrary-input Lean refinement theorem.

\subsection{Why this is relevant to machine reasoning}
The framework gives a machine explicit mathematical objects to manipulate: views, restrictions, partial realizations, typed defects, admissible corrections, and comparison witnesses. A successful step can carry both an output and a certificate of what remained fixed. A failed step can return an obstruction in a specified quotient rather than an unexplained failure flag.

The split coordinate theorem is especially direct. In its stated class, the same candidate is represented by protected information and residual data, and repair zeros only the residual coordinate. The descent theorem explains which local certificates can be assembled. The analytic results explain when a small measured defect has a quantitatively reliable meaning.

These observations motivate a machine oriented mathematical language. They do not establish an autonomous reasoning system, an empirical intelligence gain, or a general efficient algorithm. Classical choice, infinite products, quotient membership, and bounded inverse constructions can be mathematically valid without yielding executable finite procedures. Implementation claims need their own specifications and tests.

\subsection{What this monograph closes, and what it leaves open}
The book supplies a coherent account of strict realization spaces, additive and nonabelian extension problems, split tower repair, retained multiscale detail, a genuine composition-site stack, presentation transport, graph holonomy, closed-range stability, and moving exact spaces. The additional convergence and obstruction development separates small measurements from small movements, coherent coordinate families from their semantic image, and zero residual from zero transformation cost. The new bounded-budget development supplies a positive existence criterion and keeps range membership separate from regularized solvability. Its fixed-budget dual witness can certify a limit on every admissible correction, while its diagonal example computes the actual leading cost of approximation. Each layer is attached to the data that make it meaningful.

The following broader targets are not conclusions of this edition: arbitrary higher categorical coefficients and infinity-stack descent, unrestricted changes of view doctrine, norm reconstruction for every bounded infinite residual hierarchy, nonlinear or unknown-connection repair, perturbation across changing Hilbert geometries without new comparison data, and semantic essential surjectivity for unrestricted arithmetic or physical applications.

The version 1.10.0 source graph has \EditionBudgetModules\ production modules and \EditionBudgetRoots\ roots, with compiler records, declaration axiom reports, and rejection controls. All five earlier completion items KS54--KS58 retain their formal counterparts. WR01--WR07 and WR13--WR14 now have complete formal counterparts. WR09 and WR16 have formal pointwise estimates with written limiting arguments, while the remaining five WR groups retain written proofs only. The new budget results do not require finite-dimensionality, separability, injectivity, or closed range. Their robustness comparison uses one fixed correction space. Further extensions require their own hypotheses and evidence.

\begin{insight}
\textbf{The organizing principle.} Coherent realization, obstruction to extension, freedom of completion, and stability of repair belong to one mathematical conversation. Their unification is useful precisely when it does not identify them with one another.
\end{insight}

\appendix
\input{chapters/09_workshop}
\clearpage
\chap{A hypothesis atlas}\label{app:hypotheses}
A longer monograph is useful only if the reader can locate the theorem that answers the actual question. This atlas groups the main obligations by the desired conclusion. The references lead to the statements and proofs, not just to examples that resemble the intended application.

\subsection{Existence and identification}
\begin{longtable}{@{}p{.24\linewidth}p{.48\linewidth}p{.21\linewidth}@{}}
\toprule Desired conclusion & Data and hypotheses to establish & Main entry\\\midrule\endhead
A coherent family & A view system with its actual restrictions and all compatibility equations. Nonemptiness is separate. & KS01--KS03\\[4pt]
A prescribed extension & The declared partial section and an actual vanishing relative obstruction, or matching data on a transition cover. & KS04, KS09, KS14\\[4pt]
A unique extension & Existence and vanishing of the invisible coherent ambiguity, not existence alone. & KS10, KS32\\[4pt]
An infinite point from finite views & Nonempty finite fibers, or another proved lifting/compactness hypothesis. A prescribed value must lie in the persistent image. & KS54, KS58\\[4pt]
Descent of objects and witnesses & The actual groupoid data, coherent overlap isomorphisms, and a composition-covering family on the specified site. & KS34--KS36\\[4pt]
Invariance of a boundary problem & An identified boundary and the specified comparison equivalences. Literal freezing of every representative is a different construction. & KS38--KS39, KS68\\[4pt]
The same limit from two paths & Convergence of both paths in a metric space and a pairwise distance budget tending to zero. & KS74\\
\bottomrule
\end{longtable}

\subsection{Approximation and numerical meaning}
\begin{longtable}{@{}p{.24\linewidth}p{.48\linewidth}p{.21\linewidth}@{}}
\toprule Desired conclusion & Data and hypotheses to establish & Main entry\\\midrule\endhead
An exact limit of approximations & Completeness, a nonnegative lower semicontinuous residual, vanishing residuals, and summable candidate movement. & KS71--KS72\\[4pt]
A protected exact limit & In addition, continuity of the protected map and preservation along the candidate path. & KS73\\[4pt]
Exact cancellation from zero floor & Closedness of the actual admissible residual image. Without it only closure membership follows. & KS76--KS77\\[4pt]
A stable residual-to-distance bound & The required closed range or bounded inverse in the declared normed model. In a protected problem, use the restriction to $\ker p$. & KS43, KS56\\[4pt]
A controlled moving exact space & Two residual certificates, or a proved split-surjective regime with the strict smallness condition. Norm closeness alone does not preserve rank. & KS48--KS53\\[4pt]
An independent object distance & An arrow category and an extended nonnegative cost with zero identity cost and composition subadditivity. Separation is another question. & KS78--KS79\\
\bottomrule
\end{longtable}

\subsection{Bounded budgets, regularization, and solver certificates}
Every entry has a printed proof. Formal coverage is result-specific: WR01--WR07 and WR13--WR14 have complete typed counterparts, WR09 and WR16 have partial coverage, and the remaining five WR groups retain written proofs only. These statuses include the version 1.10.0 budget continuation. Each formal statement retains its own declared hypotheses.
\begin{longtable}{@{}p{.24\linewidth}p{.48\linewidth}p{.21\linewidth}@{}}
\toprule Desired conclusion & Data and hypotheses to establish & Main entry\\\midrule\endhead
An exact answer from bounded approximations & A fixed bounded linear equation between real Hilbert spaces, residual convergence, and a uniform correction bound. & WR01\\[4pt]
Strong convergence of finite repairs & Nested nonempty closed affine constraints, minimum norm answers, and a uniform energy budget. & WR02\\[4pt]
A sharp finite-budget exclusion & The correction ball, the restricted adjoint, and a dual witness in the unit ball. & WR03--WR04\\[4pt]
Robustness to model uncertainty & A fixed correction space and certified operator and data error bounds. & WR05\\[4pt]
A unique protected regularized answer & Positive regularization parameter and the actual Hilbert metric on the protected kernel. & WR06--WR07\\[4pt]
A convergent regularization limit & Actual attainment of the projected datum, not merely its membership in the closure. & WR08\\[4pt]
A noisy convergent answer & The same attained datum, vanishing parameter, and noise divided by the square root of the parameter tending to zero. & WR09, WR12\\[4pt]
A leading approximation-cost constant & The explicit diagonal model and a proof controlling its infinite Riemann-sum tails. & WR10--WR11\\[4pt]
A certified inexact optimizer & Positive parameter, the actual normal residual in the protected Hilbert norm, or an exact primal and dual gap. & WR13--WR14\\[4pt]
A basis-independent solver residual & A full independent kernel basis and its Gram inverse. The normal-equation coordinate residual is a covector. & WR15\\[4pt]
A convergent stopping schedule & Attained projected datum and separate vanishing bias, rescaled data noise, and rescaled normal residual. & WR16\\\bottomrule
\end{longtable}

\subsection{Three questions before using a finite check}
First determine whether the tested object is the actual object in the theorem. A scalar simulation can illustrate an operator-valued statement without checking every dependent fiber. Next determine whether the test checks a universal finite domain or a selected finite range. Finally identify any limit that remains outside the computation.

The new phase extension is exhaustively checked on its eight elements. The diagonal example is different: rational finite sums verify its displayed enclosures, while the proof of nonexistence uses the infinite square-summability condition. Calling both results simply ``tested'' would conceal the most important distinction.

The retained companion evidence guide and the files \path{inherited_v1.7.0/data/deepening_models.json} and \path{inherited_v1.7.0/data/written_extensions.json} retain these scopes. The general source declarations are checked in the accepted formal baseline. Printed examples and new conceptual connections do not acquire additional Lean coverage merely because they occur next to a source identifier.

\chap{Notation and source concordance}\label{app:concordance}
\subsection{Notation}
\begin{longtable}{@{}p{.25\linewidth}p{.7\linewidth}@{}}
\toprule
Symbol & Meaning\\\midrule\endhead
\(\V,F_v,f^*\) & A declared view category, its fibers, and contravariant restrictions.\\
\(\Gamma(F)\) & The strict set of coherent global families.\\
\(\Gamma_U(F),\Ext_U(s)\) & Partial sections on a full patch and their global extension fiber.\\
\(\delta,\partial,d,D\) & Respectively the view differential, tower differential, a model residual, and a Hilbert defect operator. Each target is declared where used.\\
\(p,p_U,p_L\) & A protected observable, restriction to prescribed views, and an innovation observable depending on the splitting.\\
\(\mathfrak o_U,H_U\) & The existence obstruction and the group of invisible coherent changes. Additive and nonabelian conventions are stated separately.\\
\(C_j,J_j,Q_N,S_N\) & Coarsener, chosen lift, restored coarse remainder, and retained finite reconstruction.\\
\(c_f,a_v\) & Coherent nonabelian transition labels and gauge values.\\
\(\mathfrak K_G,J_{\rm comp}\) & The cocycle patch groupoid and its composition-support topology.\\
\(\mathcal E_b(R)\) & Realizations equipped with an actual identification of their boundary with \(b\).\\
\(T_i,K_{ij},R_p\) & Transport to the root, the corresponding root loop operator, and reconstruction from the root.\\
\(\mathcal H,q,P\) & Root residual, an orthogonal kernel projector, and the explicitly defined model repair. A graph repair need not be orthogonal in the global maximum norm.\\
\(\pi_{W,U},H_U\) & The map forgetting part of the prescription and the corresponding invisible coherent ambiguity.\\
\(\phi,\operatorname{EventualHom}\) & A monotone source schedule and transport classes modulo eventual equality.\\
\(\chi_A,M=SA-BC\) & Split exact/residual coordinates and the mismatch between chosen repair lifts.\\
\(\omega,Q_n,C_n,b_n\) & A residual, its prefix-product rate, movement coefficient, and summable step budget. Here $Q_n$ is not the dyadic remainder operator $Q_N$.\\
\(\nu_B,d_W,d_W^{\mathrm{sym}}\) & Residual floor, extended directed transport distance, and its symmetric maximum.\\
\(c(x,y),\kappa(x,y)\) & The two-bit pair phase and its antisymmetric scalar-gauge invariant.\\
\(H,T,r\) & In the budget chapters, $H=\ker p$ with its inherited Hilbert metric, $T=d|_H$, and $r=df$ for the true raw residual. Measured data are declared separately.\\
\(\beta_{T,r}(M),a_\alpha\) & Minimum residual at correction budget $M$, and the quadratic regularized correction.\\
\(C,b,q\) & Closure of the correction image, its projection of the residual, and the irreducible residual floor.\\
\(v,g,e,\mathcal D_\alpha\) & Fit residual, Hilbert normal residual, optimizer error, and regularized dual objective in the solver chapters.\\
\(B,G,q\) & Kernel basis, its Gram matrix, and the coordinate normal covector. This $q$ is not the scalar residual floor used earlier.\\
\(\simeq,\simeq_L\) & A specified equivalence, and a continuous linear equivalence. Neither notation alone asserts an isometry.\\
\bottomrule
\end{longtable}

\subsection{How to read the correspondence}
The identifiers below are stable editorial identifiers, not counts of new discoveries. A row can combine several source declarations into one exposition. The named declaration, its line range, source hash, and correspondence note are recorded in \path{inherited_v1.7.0/data/claim_ledger.json} and in the retained companion Lean guide. Existence of a matching source declaration is not a mechanical proof that every word of an edited statement matches its type. The displayed hypotheses and the printed proof remain part of the review.

\textbf{F} denotes the recovered foundation tree at version 0.20.0. \textbf{M} denotes the focused moving-space tree at version 0.32.0. \textbf{S} denotes the separately recovered Hilbert holonomy source. \textbf{I} denotes a source in the canonical integrated tree, whether inherited or added during the version 1.2.0 formalization. F, M, and S retain provenance, while byte-identical copies of their referenced sources were compiled together in the accepted version 1.2.0 replay of \path{formal/KomorebiSpace/}. All \EditionClaims\ KS catalogue items retain their source correspondences. Their unchanged sources also belong to the accepted version 1.8.0 cumulative scope described below.

The sixteen WR groups are listed separately in Appendix~\ref{app:written}. Nine have complete typed counterparts, two have partial pointwise coverage, and five remain written-only. WR01 and WR02 gained their formal counterparts in version 1.8.0, followed by WR03--WR05 in version 1.10.0. Their declaration map supplements the retained KS catalogue. This does not assert that every illustrative computation or explanatory sentence is independently formalized.

\input{claim_table}

\subsection{Verification is a separate record}
The historical foundation record used a complete build followed by a partial archive replay. The separate moving-space record used two clean builds of its own closure. Neither may be reclassified as a different compilation strategy. The detailed inventories and original receipts are retained in the Lean guide and evidence archives.

The accepted version 1.2.0 record covers \EditionFormalModules\ canonical production modules and a root. Its evidence is bound by \path{inherited_v1.7.0/audit/formal/release.json}. The accepted version 1.6.0 focused record, \path{inherited_v1.7.0/audit/extension/release.json}, covers three then-new modules, three inherited dependencies, and its root. Those historical records retain their original scope.

Version 1.8.0 first rebuilt the canonical core and regularization extension together from an empty project build. All \EditionCumulativeBaseTargets\ source targets passed: \EditionCumulativeBaseModules\ production modules, the two inherited roots, and a combined root. That run audited \EditionCumulativeBaseDeclarations\ project declarations. An independent terminal check verified the complete command inventory, source and control snapshots, and compiled-object hashes.

The continuation then compiled \EditionContinuationModules\ new production modules and the updated combined root over those verified inherited objects. A fresh audit of the entire resulting environment covered \EditionCumulativeDeclarations\ declarations. That version's final source graph has \EditionCumulativeModules\ production modules and \EditionCumulativeRoots\ roots. This was a clean inherited baseline followed by an incremental extension, not a second clean build of all \EditionCumulativeTargets\ final targets. The accepted receipts and source hashes are summarized in \path{continuation/VERIFICATION_SUMMARY.json}. Appendix~\ref{app:written} lists the exact WR coverage and new public declarations.

Version 1.10.0 rebuilds the enlarged project from an empty project build. All \EditionBudgetTargets\ targets comprise \EditionBudgetModules\ production modules and \EditionBudgetRoots\ roots. The four budget modules add \EditionBudgetNewTheorems\ explicit theorem or lemma declarations, bringing that source count to \EditionBudgetTheorems. The full environment receives a fresh declaration-type and axiom audit covering \EditionBudgetDeclarations\ project declarations, including \EditionBudgetNewDeclarations\ from the new modules. All \EditionBudgetControls\ control files meet their specified outcomes: \EditionBudgetPositiveFiles\ positive files, \EditionBudgetCompilerNegativeFiles\ compiler rejection files, and \EditionBudgetPolicyNegativeFiles\ policy rejection files. The accepted run is recorded in
\begin{quote}\small\path{budget_v1.10.0/integration/runs/clean_budget_01/receipt.json}.\end{quote}
Its independent postflight checks bind the source, command inventory, dependency caches, controls, and compiled-object hashes. The precise audited declaration inventory is retained in that record. Only pinned external dependency caches are reused in this clean project replay.

These accepted runs use Lean 4.34.0-rc2 and pinned mathlib commit
\begin{quote}\small\path{aeb72f1e9981b6d35b426b16dca7759be3e7ba06}\end{quote}
The permitted standard axioms are \code{propext}, \code{Classical.choice}, and \code{Quot.sound}. The official compiler, kernel and imported caches remain trusted. Runtime restoration is not a source rebuild of mathlib. Inventory counts include auxiliary or generated declarations as specified by their audit.

Historical outlines retain their dates and scope. Later checked constructions can supersede old targets, as with the composition-site stack theorem. The current declaration, not an early roadmap, determines coverage. Printed proofs, finite tests and typed declarations remain distinct evidence.

\clearpage
\chap{Extension proofs and their verification scope}\label{app:written}
The WR catalogue is deliberately separate from the eighty KS exposition items. It began with twelve complete written theorem groups. Version 1.6.0 added WR13--WR16 and supplied a focused formal counterpart for selected parts of the catalogue. Version 1.7.0 expanded all sixteen printed proofs. The accepted version 1.8.0 continuation then formalized WR01 and WR02 and checked the enlarged cumulative source environment. Version 1.10.0 adds complete formal counterparts for WR03--WR05, covering every principal clause of the attained budget profile, sharp dual maximum, and fixed-space robustness comparison. A verified pointwise estimate still does not supply a formal proof of a separate limiting conclusion.

\begin{center}
\begin{tabularx}{\linewidth}{@{}p{.19\linewidth}X@{}}
\toprule Coverage & Exact interpretation\\\midrule
Complete: \EditionWRComplete & WR01--WR07, WR13 and WR14. The principal universal statements have typed counterparts in the accepted cumulative scope. Explanatory examples retain their stated separate scope.\\[5pt]
Partial: \EditionWRPartial & WR09 and WR16. The fixed-parameter and pointwise solver/data inequalities are formal. Their parameter-limit consequences still use written WR08.\\[5pt]
Written only: \EditionWRWrittenOnly & WR08, WR10--WR12 and WR15. Full proofs are printed. No complete typed counterpart is assigned to these groups.\\\bottomrule
\end{tabularx}
\end{center}

The version 1.8.0 verification first compiled \EditionCumulativeBaseTargets\ targets from an empty project build, combining the inherited core and regularization extension. It then compiled the \EditionContinuationModules\ new WR01--WR02 modules and updated combined root over those verified objects. The fresh full audit covered \EditionCumulativeDeclarations\ declarations in the resulting environment. All \EditionCumulativePositiveFiles\ positive regression files, \EditionCumulativeCompilerNegativeFiles\ compiler rejection files, and \EditionCumulativePolicyNegativeFiles\ policy rejection files passed their specified checks. The final declaration inventory includes compiler-generated and auxiliary items. It is not a count of independent mathematical results.

The version 1.10.0 budget continuation instead recompiles all \EditionBudgetTargets\ project targets from an empty project build. Its \EditionBudgetModules\ production modules include four new modules with \EditionBudgetNewTheorems\ explicit theorem or lemma declarations. A fresh full-environment audit and \EditionBudgetControls\ control outcomes accompany that replay. The final receipt and independent postflight record are under \path{budget_v1.10.0/integration/runs/clean_budget_01/}. This is new formal verification, while the earlier numerical suites remain historical records.

The historical source files retain their original bytes. The final cumulative sources are supplied in \path{formal/}, and the accepted continuation records in \path{continuation/}. The compiler, kernel, and pinned imported dependency caches remain trusted. The source-type comparison is an editorial audit, not automatic verification of natural-language proofs. Appendix~\ref{app:concordance} distinguishes the clean baseline from the incremental continuation.

\subsection{Existence and budget geometry}
\textbf{WR01. Bounded approximate feasibility.} The real Hilbert operator and residual datum are fixed. A uniform correction-norm bound gives an exact preimage within the same bound. The printed proof uses a weakly convergent subsequence. The formal proof uses weak-star compactness of a dual ball and the Riesz isometric equivalence, with no separability or closed-range assumption. Its consequences include the range-membership equivalence and divergence of every approximating sequence's correction norms when the datum is unattained. The protected bridge applies the theorem to the actual kernel of the observable and constructs an exact repaired candidate within the original correction budget.

\textbf{WR02. Nested affine minimum corrections.} Nonempty closed affine constraint sets are nested in one real Hilbert space. The formal development constructs their minimum-norm points, proves the exact energy identity and norm monotonicity, and establishes the equivalence between a nonempty intersection and a uniform bound on the minimum norms. Under that bound, it proves strong convergence to the unique minimum-norm point of the intersection and the limiting energy identity. If the intersection is empty, the minimum norms tend to infinity. The affine hypothesis matters: the more general convex argument supplies an energy inequality rather than equality.

\textbf{WR03. Attained budget profile.} A finite nonnegative budget has an attained residual minimum even when the full operator range is not closed. The formal proof derives closedness of the budget image from WR01 and applies Hilbert projection. It proves monotonicity, convexity, the Lipschitz bound, the unlimited-budget floor, and both finite-budget equivalences. The printed proof retains its weak topology argument.

\textbf{WR04. Sharp budget support.} The general support formula uses the actual restricted adjoint. The formal result now proves the full maximum formula at every nonnegative budget. A minimizing correction supplies a normalized residual witness when the residual is nonzero, and the zero witness handles the zero-residual case. Both the universal lower certificate and its attainment are proved.

\textbf{WR05. Robust budget comparison.} The operators act on the same correction space, with operator and data errors bounded by $\eta$ and $\delta$. The formal result includes the absolute profile bound, candidate-level estimates, the nominal dual lower certificate, and strict robust exclusion. The bound $\delta+\eta M$ does not apply automatically when the protected kernel changes.

\subsection{Regularization and its limit}
\textbf{WR06. Constructed regularized minimizer.} Positive $\alpha$ makes $T^*T+\alpha I$ coercive. The formal proof constructs an actual continuous linear equivalence, defines the correction using its inverse, and proves the energy identity and uniqueness. The manuscript's self-adjoint closed-range proof and the formal coercive-form proof establish the same stated result. The bridge uses the existing \path{ProtectedHilbertRepair.restricted} operator.

\textbf{WR07. Parameter and budget comparison.} The formal statements compare two positive parameters, prove the two norm monotonicities, and show that no correction of smaller or equal norm has a lower residual than the regularized answer. The normalized residual has unit norm and attains the support expression whenever it is nonzero. The Lean formulation gives the optimizer and universal competitor inequality directly rather than introducing a second numerical profile definition.

\textbf{WR08. The regularization limit.} Projection onto the closure of the range is separated from actual range membership. Strong convergence of corrections requires the projected datum to be attained. Otherwise the correction norms diverge. These infinite-limit arguments remain written-only.

\textbf{WR09. Data sensitivity and parameter choice.} The focused proof supplies the coefficients $1/(2\sqrt\alpha)$ for correction noise and one for fitted-data noise. The scalar sharpness calculations and the parameter-choice convergence consequence retain their printed proofs. Formalizing a fixed-parameter inequality does not formalize its convergence application automatically.

\textbf{WR10--WR12. The explicit diagonal models.} The exact regularized family, omitted-tail bounds, matched-budget limit $\pi/4$, and attainable noisy datum remain written with exact finite enclosure tests. No general infinite-sum formalization or new diagonal Lean instance is claimed.

\subsection{Inexact solution and coordinate interpretation}
\textbf{WR13. A posteriori solver bounds.} The correction error, fitted error, and objective excess are bounded in terms of the actual normal residual and positive parameter. The focused theorem quantifies over every proposal in the Hilbert source. The protected bridge transfers its candidate-distance bound to the actual kernel subtype. Scalar sharpness examples are separated from the principal theorem and checked independently.

\textbf{WR14. A regularized primal and dual certificate.} An exact two-square identity proves weak duality, a dual maximizer constructed from the primal correction, global objective bracketing, and a correction-distance bound. The witness is arbitrary. Setting it to the proposal's fit residual gives the exact normal-residual gap.

\textbf{WR15. The coordinate covector metric.} With orthonormal ambient and residual coordinates, a full independent kernel basis gives $g=BG^{-1}q$ and $\norm g^2=q^{\mathsf T}G^{-1}q$. The invariance proof and its invertibility hypotheses are printed and rationally tested. They have not been compiled as a new matrix theorem.

\textbf{WR16. A justified stopping rule.} The focused theorem proves the pointwise sum of the solver and data bounds. The printed convergence consequence uses written WR08 and the additional assumptions $\alpha\to0$, $\delta/\sqrt\alpha\to0$, and $\epsilon/\alpha\to0$. The displayed diagonal schedule also uses written WR12.

\subsection{Public declarations for WR01 and WR02}\label{app:wr-source-map}
The following map identifies the principal clauses without reproducing the full Lean types. Those types, their actual hypotheses, and their audited axiom dependencies are recorded in \path{continuation/FORMAL_COVERAGE.json}. Every declaration below is in the accepted cumulative environment. This map does not assign separate typed coverage to every worked example or background proof.

For WR01, the declarations belong to the namespace
\begin{quote}\small\path{KomorebiSpace.BoundedApproximateRealization}.\end{quote}
\begin{longtable}{@{}p{.30\linewidth}p{.64\linewidth}@{}}
\toprule Printed clause & Declaration\\\midrule\endhead
Exact solution within the same budget & \path{exists_solution_of_bounded_residual_tendsto_zero}\\[4pt]
Range-membership equivalence & \path{mem_range_iff_bounded_approximation}\\[4pt]
Escape for unattained data & \path{norm_tendsto_atTop_of_not_mem_range}\\
\bottomrule\end{longtable}

These declarations concern bounded linear maps between real Hilbert spaces. Residual convergence is stated either as convergence of the residual norm to zero or, equivalently, convergence of the images to the fixed datum. The conclusion preserves the supplied norm bound. It does not assert strong convergence of the original approximating sequence.

\Needspace{22\baselineskip}
For WR02, the declarations belong to the namespace
\begin{quote}\small\path{KomorebiSpace.NestedMinimumCorrections}.\end{quote}
\begin{longtable}{@{}p{.30\linewidth}p{.64\linewidth}@{}}
\toprule Printed clause & Declaration\\\midrule\endhead
Constructed minimum corrections & \path{affineCorrections_spec}\\[4pt]
Nondecreasing minimum norms & \path{nested_affine_norm_monotone}\\[4pt]
Finite-stage energy equality & \path{nested_affine_energy}\\[4pt]
Intersection and boundedness equivalence & \path{nested_affine_intersection_iff_bounded}\\[4pt]
Strong limit, unique minimum, and limiting energy equality & \path{nested_affine_strong_realization}\\[4pt]
Escape when the intersection is empty & \path{nested_affine_escape}\\
\bottomrule\end{longtable}

The Lean family is indexed by the natural numbers and is antitone under inclusion. Each affine set is explicitly nonempty and closed. Completeness is assumed for the common real Hilbert space. In the convergence theorem, the boundedness assumption concerns the constructed minimum-norm corrections themselves.

The protected application belongs to
\begin{quote}\small\path{KomorebiSpace.ProtectedBoundedRealization},\end{quote}
with public declaration
\begin{quote}\small\path{exists_exact_repair_of_bounded_approximation}.\end{quote}
It uses corrections in the actual subtype $\ker p$ with its inherited norm. For a raw candidate $f$, it constructs $y$ satisfying $Dy=0$, $py=pf$, and $\norm{f-y}\le M$. The defect's source and target are real Hilbert spaces. The observable's codomain need only be a real normed space. No closed-range hypothesis is introduced.

\subsection{Public declarations for WR03--WR05}
The budget continuation adds four modules. All three theorem groups use bounded linear maps between real complete Hilbert spaces. No finite-dimensionality, separability, injectivity, or closed-range assumption is added. The feasible correction ball is the same in the two problems compared by WR05, and zero budget is included.

\begin{samepage}
The WR03 existence and floor declarations belong to
\begin{quote}\small\path{KomorebiSpace.BudgetProfile}.\end{quote}
\begin{center}\begin{tabular}{@{}p{.30\linewidth}p{.64\linewidth}@{}}
\toprule Printed clause & Declaration\\\midrule
Closed budget image & \path{isClosed_budget_image}\\[4pt]
Attainment at each budget & \path{exists_budget_minimizer}\\[4pt]
Printed infimum interpretation & \path{beta_eq_iInf}\\[4pt]
Limit at infinite budget & \path{beta_tendsto_floor}\\[4pt]
Finite-budget equality to the floor & \path{beta_eq_floor_iff}\\[4pt]
Exact feasibility at zero residual & \path{beta_eq_zero_iff}\\
\bottomrule\end{tabular}\end{center}
\end{samepage}

\begin{samepage}
The remaining profile geometry belongs to
\begin{quote}\small\path{KomorebiSpace.BudgetProfileGeometry}.\end{quote}
\begin{center}\begin{tabular}{@{}p{.30\linewidth}p{.64\linewidth}@{}}
\toprule Printed clause & Declaration\\\midrule
Nonincreasing profile & \path{beta_antitone}\\[4pt]
Zero budget & \path{beta_zero}\\[4pt]
Convexity on nonnegative budgets & \path{beta_convex_combination}\\[4pt]
Two-sided profile difference bound & \path{beta_profile_bounds}\\
\bottomrule\end{tabular}\end{center}
\end{samepage}

\begin{samepage}
For WR04, the declarations belong to
\begin{quote}\small\path{KomorebiSpace.BudgetDuality}.\end{quote}
\begin{center}\begin{tabular}{@{}p{.30\linewidth}p{.64\linewidth}@{}}
\toprule Printed clause & Declaration\\\midrule
Universal dual lower bound & \path{dualValue_le_beta}\\[4pt]
Attaining unit-ball witness & \path{exists_attaining_dual_witness}\\[4pt]
Literal attained maximum & \path{sharp_budget_duality}\\[4pt]
Strict exclusion certificate & \path{strict_exclusion_certificate}\\
\bottomrule\end{tabular}\end{center}
\end{samepage}

The maximum is expressed as \code{IsGreatest} for the set of dual values over the closed unit ball. This includes an actual attaining witness. It does not infer attainment merely from a supremum or exchange an infimum with a supremum.

\begin{samepage}
For WR05, the declarations belong to
\begin{quote}\small\path{KomorebiSpace.BudgetRobustness}.\end{quote}
\begin{center}\begin{tabular}{@{}p{.30\linewidth}p{.64\linewidth}@{}}
\toprule Printed clause & Declaration\\\midrule
Candidate uncertainty estimate & \path{candidate_residual_le}\\[4pt]
Absolute profile comparison & \path{beta_abs_sub_le}\\[4pt]
Nominal dual lower certificate & \path{nominal_dual_lower_certificate}\\[4pt]
Strict robust exclusion & \path{excludes_budget_of_nominal_dual_certificate}\\
\bottomrule\end{tabular}\end{center}
\end{samepage}

The operator- and datum-error bounds themselves force $\eta,\delta\ge0$, so the formal profile comparison does not repeat those sign assumptions as separate arguments. Its budget is explicitly nonnegative. The unrestricted candidate-level estimate does not require completeness.

\subsection{Exact implementation and remaining formalization}
The retained version 1.7.0 regularization suite records \EditionRegularizationChecks\ assertions and its inexact-certificate suite \EditionInexactChecks. These are historical numerical records, not new runs performed for this edition. The independent oracle solves ambient constrained systems by determinant formulas rather than using the certificate program's kernel coordinates. The verifier checks a full kernel basis, retains its Gram metric, preserves the observable, and distinguishes measured fit from the actual residual. All outputs are exact rationals or squared bounds. It recomputes records and rejects altered claims without calling an exact optimizer. These finite tests are not universal proofs that the Python programs implement the Lean definitions.

The next central formalization target is WR08, the regularization limit, which supports the remaining limiting conclusions of WR09 and WR16. WR10--WR12 and WR15 also retain their written status. The completed budget results do not automatically discharge those separate obligations.


\input{chapters/12_proof_tools}
\input{chapters/08_references}
\end{document}
